\begin{prop}\label{key-prop}
    Let $\beta \in \fS$ be a basic $\fS$-set.
    \begin{enumerate}
        \item $D_\beta$ is an $\fS$-set, and $\wt(D_\beta \setminus \beta) < \wt(\beta)$.
        \item For every basic set $\alpha \in \fS$ there is a constant $N^\beta_\alpha$ such that
        \[
            \sum_{b : a \leq b \in \beta^*} \binom{(k)^d - a}{(k)^d - b} = N^\beta_\alpha \qquad (a \in \alpha).
        \]
        \item Every element of $\beta$ is dominated by a unique element of $\beta^*$.
        \item Either $\wt(\beta) = \wt(D_\beta \setminus \beta) + 1$ or $\beta^* \subset \{0, k\}^d$.
    \end{enumerate}
\end{prop}
\begin{proof}
    Let $w = \wt(\beta)$ and $D = D_\beta$.
    Then, since the polynomials $p^a_{b,c}(m)$ have positive leading coefficient whenever they are nonzero,
    the degree of $p^a_{\beta,\beta}(m)$ is, by the previous lemma,
    \[
        \deg(p^a_{\beta,\beta}(m)) = \max_{b, c \in \beta} \deg(p^a_{b,c}(m)) \le \max_{b, c \in \beta} \wt(\min(b, c)) \le w,
    \]
    with equality holding if and only if $a \leq b = c \in \beta^*$, i.e., if and only if $a \in D$.
    Since~$p^a_{\beta,\beta}(m)$ should depend only on the cell of $\fS$ containing $a$,
    it follows that $D$ is an $\fS$-set.
    Since $\beta$ is basic, $\beta \subset D$, and since $\beta \supset \beta^*$ we have $\wt(D \setminus \beta) < \wt(\beta)$.
    
    Moreover if $a \in D$ then the leading term of $p^a_{\beta,\beta}$ is
    \[
        \lambda(p^a_{\beta,\beta}(m)) = \sum_{b : a \leq b \in \beta^*} \binom{(k)^d - a}{(k)^d - b} \frac1{b!} m^w,
    \]
    and again this should depend only on the cell $\alpha$ containing $a$.
    {Taking $\alpha = \beta$ and $a\in \beta^*$, it follows that $a!$ is a constant for $a \in \beta^*$.}
    Hence (2) holds (if $\alpha$ is not contained in $D$ then $N^\beta_\alpha = 0$).
    
    Next apply (2) with $\alpha = \beta$.
    Taking $a \in \beta^*$ shows $N^\beta_\beta = 1$, so (3) holds.
    
    Finally let $b \in \beta^*$ and suppose $a = b - e_i \ge 0$.
    If $a \in \beta$ then we get
    \[
        \binom{k-b_i+1}{k-b_i} = N^\beta_\beta = 1,
    \]
    so $b_i = k$.
    Hence either $\wt(b) = \wt(\beta)+1$ or $b_i \in \{0, k\}$ for all $i$, which implies (4).
\end{proof}
